\chapter{Time alignment and delays}
\label{ch:timing}

Identification pairs an input (thrust) with an output (motion) sample by sample. If the two
are shifted in time, the estimator will explain the shift with whatever coefficient fits it
best, and that is usually the inertia.

\section{Clocks in a rosbag}

Every message in an MCAP file has a \code{log_time}, the moment the recorder received it on
the Jetson. Messages with a header also carry \code{header.stamp}, set by the publisher.
\begin{itemize}
  \item The EKF (\code{/odometry/filtered/local}) is used on its \code{header.stamp}.
  \item \code{/pixhawk/servo_output_raw} is a bare \code{Int16MultiArray} with \emph{no header},
        so its \code{log_time} is the only time available. That is when the Jetson
        recorded it, not when the Pixhawk applied it.
  \item The BMS voltage on \code{/diagnostics} is also taken on \code{log_time}.
\end{itemize}
Pairing servo \code{log_time} with EKF \code{header.stamp} is meaningful only if the EKF stamps
its messages close to when they are recorded. \code{verify_timestamps()} in \file{bag_io.py}
therefore reports, before any propagation: the rates and gaps of each stream, backwards or
duplicate timestamps, the median and 95th percentile of \code{log_time - header.stamp} for
the EKF (a warning above \qty{1}{\second}), and the overlap window. Problems are reported
and never silently corrected.

\section{Two different delays}

\Cref{tab:rates} lists the relevant streams. Sample rate is not the bottleneck: the EKF has
about 44 samples per \qty{1.4}{\second}. The delays are the bottleneck.

\begin{table}[htbp]
\centering
\caption{Sample periods of the streams used for identification, measured on
\code{aut_no_z_surface_big_40_grid_03}.}
\label{tab:rates}
\begin{tabular}{@{}l S[table-format=3.0] l@{}}
\toprule
Stream & {Period / \unit{\milli\second}} & Note \\
\midrule
EKF \code{/odometry/filtered/local} & 33 & \qty{30}{\hertz} \\
DVL \code{/sensors/dvl/odometry_cov} & 132 & \qty{7.6}{\hertz} \\
DVL \code{/sensors/dvl/velocity} & 230 & \qty{4.3}{\hertz} \\
\code{servo_output_raw} & 100 & \qty{10}{\hertz} (\qty{2}{\hertz} in older rectangle bags) \\
IMU \code{/imu/data} & 12 & about \qty{86}{\hertz} \\
\bottomrule
\end{tabular}
\end{table}
\FloatBarrier

\paragraph{1. Actuation and transport delay.} This is the time from the PWM being logged to
the vehicle actually accelerating: MAVLink, bridge and recorder transport (tens of
\unit{\milli\second}) plus motor and propeller spin-up (\qtyrange{100}{300}{\milli\second}). Transport and
spin-up cannot be separated without a Pixhawk-side timestamp, and they do not need to be.
\code{thrust.servo_delay_s} shifts the servo \code{log_time} by the total. It is currently \num{0.0}.

\begin{revised}[title=Revised 14 October 2026: the sign of \code{servo_delay_s}]
The first edition said that \code{servo_delay_s} shifts the log time ``by the total, which is
exactly the quantity needed''. That is true in magnitude but hides the sign. The scripts
\emph{subtract} it from the servo \code{log_time}, so a positive value makes the thrust act
\emph{before} its log entry. Spin-up makes the thrust act \emph{after} it, so a measured lag
$L$ (IMU edge minus logged PWM edge) is entered as $\code{servo_delay_s} = -L$. The synthetic
test bench confirms this numerically through two independent code paths (\cref{tab:synth-prop,tab:synth-comp}).
\end{revised}

\paragraph{2. Filter lag.} The EKF smooths its inputs. The DVL and the EKF observe the same
motion, so the delay between them isolates the EKF lag, while the actuation delay is common
to both and cancels.

\section{How a delay biases the inertia}

A response that appears \emph{later} than it really happened looks \emph{slower}, and a slower
response is explained by a \emph{heavier} vehicle. The bias always overestimates $\mtot$, and
therefore $\Xud$. Drag estimates are hardly affected, because they come from steady states or
from decays lasting many seconds.

A simple model makes this concrete. Suppose $\mtot$ is estimated as $T$ divided by the slope
of a straight line fitted to $u(t)$ over a window $[0, W]$ after the logged thrust onset. If
the true motion starts $d$ seconds late, the velocity ramp occupies only $W - d$ of the
window, and for short windows $\hat m \approx \mtot\,W/(W-d)$. A second error appears even
without delay: over a window comparable to $\taueff$, drag already bends the curve, and the
fitted slope is lower than $T/\mtot$.

\begin{figure}[htbp]
\centering
\begin{subfigure}{0.48\textwidth}
\centering
\begin{tikzpicture}
\begin{axis}[notes,xlabel={$t$ / \unit{\second}},ylabel={$u$ / \unit{\meter\per\second}},xmin=0,xmax=1.5,ymin=0]
\addplot[black] table[x=t,y=true]{data/delay.dat};
\addplot[red,dashed] table[x=t,y=delayed]{data/delay.dat};
\legend{true,delayed \qty{0.2}{\second}}
\end{axis}
\end{tikzpicture}
\caption{Step response at \qty{20}{\newton}}
\end{subfigure}\hfill
\begin{subfigure}{0.48\textwidth}
\centering
\begin{tikzpicture}
\begin{axis}[notes,xlabel={Fit window $W$ / \unit{\second}},ylabel={$\hat m / \mtot$},xmin=0.25,xmax=1,ymin=0,ymax=6]
\addplot[black] table[x=W,y=nodelay]{data/delay_bias.dat};
\addplot[red,dashed] table[x=W,y=delayed]{data/delay_bias.dat};
\addplot[gray,dotted,forget plot] coordinates {(0.25,1) (1,1)};
\legend{no delay,delay \qty{0.2}{\second}}
\end{axis}
\end{tikzpicture}
\caption{Inertia from a straight-line slope fit}
\end{subfigure}
\caption{Simulated surge step response at \qty{20}{\newton} with and without a
\qty{0.2}{\second} delay, and the inertia estimate $\hat m = T/\text{slope}$ from a linear fit
over $[0, W]$ normalised by the true $\mtot$, for the dummy BlueROV2 coefficients.}
\label{fig:delay}
\end{figure}
\FloatBarrier

\Cref{fig:delay} shows both effects for the dummy coefficients. Because $\taueff$ at
\qty{20}{\newton} is only \qty{0.19}{\second} (\cref{tab:taueff}), no window gives a good
answer. Short windows are dominated by the delay ($\hat m/\mtot$ above 3 at
$W = \qty{0.3}{\second}$), and long windows by the drag ($\approx1.7$ even without delay
at $W = \qty{0.5}{\second}$).

\begin{keyidea}
Reading the inertia off a slope is fragile when $\taueff$ is short. Two remedies, which can be combined:
\begin{enumerate}[nosep]
  \item \textbf{Measure the delay first} (next section) and set \code{servo_delay_s}.
  \item \textbf{Fit the whole step as a trajectory}: integrate \eqref{eq:surge} with the
        recorded thrust and adjust $\mtot$, $X_u$, $\Xuu$ (and optionally the delay) until the
        simulated speed matches the measurement (output error, \cref{sec:output-error}). Drag
        is then modelled rather than ignored, so the window problem disappears.
\end{enumerate}
\end{keyidea}

\section{Measuring the actuation delay with the IMU}

A thrust step is a near-instantaneous step in \emph{acceleration}, but only the start of a
ramp in \emph{velocity}. The start of a ramp is hard to locate in noise, while an edge is easy.
The IMU's \code{linear_acceleration.x} is therefore the right instrument for timing. It is
also sampled about ten times faster than the DVL. The signal is strong: about
\qty{17}{\newton} on about \qty{20}{\kilogram} is \qty{0.85}{\meter\per\second\squared}, far
above MEMS noise. The \code{thrust_cut} edge is just as sharp, so every run yields two delay
measurements.

\paragraph{Specific force.} An accelerometer measures specific force
$\vect{f} = \vect{a} - \vect{g}^b$, not acceleration. In FLU/ENU,
$\vect{g}^b = \vect{R}^\top[0,0,-g]^\top$, whose $x$ component for a pure pitch is
$+g\sin\theta$ (nose-down pitch tilts gravity forward). Hence
\begin{equation}
  \dot u \approx a_x = f_x + g\sin\theta .
\end{equation}
The trim pitch adds a constant that must be removed before using the magnitude. The pitching
caused by \code{MOT_1} sitting above the CoM arrives \emph{after} the edge, so it does not
affect the edge timing.

\paragraph{A bonus.} The accelerometer also gives the inertia without differentiating any velocity:
$\mtot = T/a_x$ at the instant of the step, while the damping is still close to zero.

\section{Thruster delay: what is measured, and telling it apart from mixer errors}
\label{sec:thruster-delay}

\begin{revised}[title=Added 30 October 2026,breakable]
\paragraph{What the IMU onset contains.} Acceleration responds to force without lag
($a = F/m$); the inertia only slows the velocity. After a step from rest, where the drag is still
close to zero, the shape of the measured acceleration is therefore the shape of the rising
thrust. The onset measured against the logged PWM is the sum of the logging and transport delay,
the ESC response and the motor and propeller spin-up. That sum is the quantity the fit needs,
because the fit pairs the \emph{logged} PWM with the motion, so it does not have to be split.
Spin-up is a ramp rather than a pure delay, so a delay plus a first-order lag is fitted to the
acceleration after each edge: a lag cannot be removed by shifting time stamps. Spin-up,
spin-down and reversal through neutral are measured separately, since their timings differ.

\paragraph{Isolating the parts.} Where the parts are wanted individually:
\Cref{tab:delay-parts} lists the measurement for each. Per thruster, the bench is the clean
route, since one unit is mounted at a time. On the vehicle the mixer commands axes, not motors;
ArduSub's motor test (\code{MAV_CMD_DO_MOTOR_TEST}) drives a single motor, subject to the
restrictions of the installed firmware. Since the units of one model share the same ESC, testing
two or three units on the bench is enough; the model uses one delay per axis group.

\paragraph{Timing mismatch versus static error.} A pure yaw command drives the front and rear
lateral thrusters (\code{MOT_6}, \code{MOT_2}) in opposite directions, with mixer factors chosen
so that their sway forces cancel. A mismatch shows up as sway the command should not produce, and
its time signature tells the cause:
\begin{itemize}
  \item a \textbf{timing mismatch} gives a short lateral kick at the start of the step, an
        opposite kick at its end, and nothing during the hold;
  \item a \textbf{static error} (mixer ratio, lever arm or unequal thrust) gives a steady lateral
        force for the whole hold, roughly proportional to the command.
\end{itemize}
A step held for several seconds shows both at once. The steps start from rest, so that no drag
and no speed-dependent coupling are present, and the lateral acceleration is corrected for the
IMU's offset from the CoM (about \qty{5}{\centi\meter} forward, which adds $\dot r$ times that
offset in a turn). Pure sway steps cross-check the result: a short yaw-rate kick at their edges
indicates the same mismatch. Steps in both yaw directions distinguish a slow \emph{unit} (the
same thruster always leads) from a slow \emph{reverse branch} (the leader changes with the
direction, because one thruster always runs in reverse). With two thrusters and two axes, the
allocation can also be inverted to reconstruct each thruster's force from the lateral and yaw
accelerations; unknown inertias scale these curves but do not shift their onsets.

\paragraph{Which static error.} The logged \code{servo_output_raw} shows what the mixer actually
sent to each motor, so the mixer's arithmetic is checked from the log alone. With the PWM correct,
a remaining steady mismatch has three possible causes that look identical on the vehicle: the
geometry assumed by the mixer differs from the real lever arms, the two units deliver different
thrust at the same PWM, or the reverse branch is weaker. The bench curves of each unit in both
directions remove the last two, and whatever steady mismatch then remains measures the geometry.

\paragraph{Refining the allocation and the mixer.} Once the hardware is final and the bench
curves exist, these steady cross-axis tests give the actual lever ratios. Two things can then be
corrected, for different reasons. The allocation in \file{model.yaml} is what the model uses; the
model takes the logged PWM, so its correctness does not depend on the mixer, only on the
allocation. The mixer factors in the ArduSub custom frame (firmware fork) determine how well a
pilot or controller command stays on one axis; correcting them improves the vehicle's handling
and requires a firmware rebuild.
\end{revised}

\begin{table}[htbp]
\centering
\caption{Parts of the delay between the logged PWM and the force on the hull, and how each can be
measured on its own.}
\label{tab:delay-parts}
\begin{tabular}{p{0.3\linewidth}p{0.62\linewidth}}
\toprule
Part & Measurement \\
\midrule
Logging and transport & Publish time of \code{manual_control} against the matching change in
\code{servo_output_raw}, or the Pixhawk's own \code{time_usec} (currently dropped by the bridge) \\
ESC and thruster & Thrust bench with a PWM step, load cell read at full rate (NAU7802: up to
\qty{320}{\hertz}) \\
Motor spin-up alone & Onset of the motor vibration in the raw high-rate IMU signal, or RPM from
ESC telemetry if the ESC provides it \\
Total, per axis group & IMU acceleration after a step from rest (\cref{sec:thruster-delay}) \\
\bottomrule
\end{tabular}
\end{table}
\FloatBarrier

\section{Measuring the filter lag}

On the same step, the DVL velocity and the EKF velocity are compared. Three outcomes are possible:
\begin{itemize}[nosep]
  \item if they rise together, the EKF is used for everything;
  \item if the EKF visibly trails, inertia is fitted against the DVL and the EKF is kept for slower quantities;
  \item if both look smeared, the manoeuvre is too slow for the sample rates, and the
        remedy is a sharper command.
\end{itemize}
The DVL is \emph{less} lagged than the EKF, but it is not free of lag. The A50's internal ping
averaging should be checked before it is treated as the reference.

\begin{questions}
\begin{enumerate}
  \item Why does a delay make the vehicle look heavier, but leave the drag almost unchanged?
  \item Why is the IMU better than the DVL for timing the thrust onset?
  \item The accelerometer reads $f_x = \qty{-0.17}{\meter\per\second\squared}$ while the vehicle
        is at rest with $\theta_0 = \qty{1}{\degree}$. Is that consistent?
\end{enumerate}
\end{questions}
